% !TEX program = xelatex
% =====================================================================
%  BÁO CÁO ĐỒ ÁN TRÍ TUỆ NHÂN TẠO, HỌC PHẦN AIP202
%  Khoa Công nghệ Thông tin, Trường Đại học Kiến trúc Đà Nẵng
%
%  Build:  latexmk -xelatex do-an.tex
%  Overleaf: Menu > Compiler > XeLaTeX, rồi bấm Recompile.
%
%  Khung tài liệu chuyển thể từ mẫu đồ án của Trường Đại học Nha Trang
%  (github.com/nd-hung/thesis-template, tác giả TS. Nguyễn Đình Hưng).
%  Toàn bộ nhận diện, phông chữ, bảng màu, trang bìa và cấu trúc chương
%  được thay theo chuẩn của DAU và theo đề cương học phần AIP202.
% =====================================================================
% \documentclass phải nằm ở file này, không được chuyển vào settings.tex,
% vì Overleaf dò file chính bằng cách tìm file có chứa \documentclass.
\documentclass[12pt, a4paper, oneside]{book}
% !TEX root = ../do-an.tex
% =====================================================================
%  styles/settings.tex
%  Các thiết lập chung, nạp từ file chính do-an.tex.
%  SINH VIÊN KHÔNG SỬA FILE NÀY.
%
%  Lưu ý: lệnh khai báo lớp tài liệu nằm ở file do-an.tex chứ không nằm
%  ở đây. Lý do: Overleaf dò file chính bằng cách tìm chuỗi khai báo lớp
%  tài liệu trong các file. Nếu chuỗi đó xuất hiện ở file này, Overleaf
%  sẽ biên dịch nhầm file này và báo lỗi "no legal end found", vì file
%  này chỉ chứa thiết lập chứ không chứa phần thân tài liệu.
% =====================================================================

% Toàn bộ nhận diện DAU: phông, màu, bìa, hộp trình bày, header
\usepackage{styles/dau}
\dauLayout

% Giãn dòng và khoảng cách đoạn theo quy cách đồ án
\linespread{1.25}
\usepackage[skip=8pt plus 1pt, indent=15pt]{parskip}

% Ghi chú: mẫu gốc dùng gói atbegshi để vứt trang đầu, vì trang bìa của
% mẫu gốc dựng bằng minipage nên sinh ra một trang trắng. Trang bìa ở
% đây dựng bằng môi trường titlepage nên không có trang trắng, và nếu
% giữ thủ thuật đó thì chính trang bìa sẽ bị vứt mất.


% =====================================================================
%  PHẦN EM ĐIỀN: thông tin đồ án
% =====================================================================
% Loại báo cáo
\crname{BÁO CÁO ĐỒ ÁN TRÍ TUỆ NHÂN TẠO}
% Học phần
\cmonhoc{Học phần AIP20201 \textperiodcentered{} 2 tín chỉ}
% Tên đề tài. Dùng \\ để ngắt dòng cho cân đối trên bìa.
\ctname{Xây Dựng Hệ Thống Hỗ Trợ Tìm Phòng Trọ Đà Nẵng}
% Họ tên sinh viên và mã số sinh viên
\cstuname{Trần Nam Phương}
\cmssv{2351220238}
% Lớp
\clopsv{23CT3}
% Ngành
\csCouncil{Công Nghệ Thông Tin}
% Giảng viên hướng dẫn
\csSupervise{TS. Đỗ Phúc Hảo}
% Năm học
\cttime{2026--2027}
% Khoa quản lý
\csdeptname{KHOA CÔNG NGHỆ THÔNG TIN}
% Nơi thực hiện
\csplace{Đà Nẵng}

\begin{document}

% Bìa trang trọng có khung kép. Muốn dùng bìa hiện đại theo bộ giáo
% trình của khoa thì thay dòng dưới bằng \coverpagehiendai
\coverpage

% ---------------------------------------------------------------------
%  PHẦN ĐẦU: đánh số trang bằng số La Mã thường
% ---------------------------------------------------------------------
\frontmatter
\dauPageStyle

% ---------------------------------------------------------------------
%  LỜI CAM ĐOAN
% ---------------------------------------------------------------------
\begin{declaration}

Tôi xin cam đoan đồ án này là công trình do chính tôi thực hiện dưới sự
hướng dẫn của \dauGVHD. Các nội dung nghiên cứu, số liệu và kết quả trình
bày trong đồ án là trung thực và chưa từng được công bố trong bất kỳ công
trình nào khác.

Mọi số liệu, hình ảnh, mã nguồn và kết quả của tác giả khác được sử dụng
trong đồ án đều được trích dẫn nguồn gốc rõ ràng trong phần tài liệu tham
khảo. Toàn bộ phần mềm và thư viện được dùng đều có giấy phép cho phép sử
dụng. Những phần có sử dụng công cụ trí tuệ nhân tạo hỗ trợ đã được nêu rõ
trong báo cáo.

Nếu phát hiện có bất kỳ sự gian lận nào, tôi xin chịu hoàn toàn trách nhiệm
trước Khoa và Nhà trường.

\vspace{1.4em}
\hfill\okyten{\dauSinhVien}{Sinh viên thực hiện}

\end{declaration}

% ---------------------------------------------------------------------
%  LỜI CẢM ƠN
% ---------------------------------------------------------------------
\begin{acknowledgments}

Em xin gửi lời cảm ơn chân thành đến \dauGVHD, người đã tận tình định hướng
đề tài, góp ý về cách thiết kế hệ gợi ý và theo sát tiến độ của em trong
suốt quá trình thực hiện đồ án. Em cũng xin cảm ơn quý thầy cô Khoa Công
nghệ Thông tin, Trường Đại học Kiến trúc Đà Nẵng đã trang bị cho em nền
tảng kiến thức để hoàn thành công việc này. Em xin cảm ơn các bạn sinh viên
và một số chủ nhà trọ tại khu vực Đà Nẵng đã nhiệt tình cung cấp thông tin
phòng trọ và tham gia khảo sát, giúp em có dữ liệu thực tế để xây dựng và
đánh giá hệ thống. Cuối cùng, em xin cảm ơn gia đình và bạn bè đã luôn động
viên em trong thời gian học tập.

\vspace{1.4em}
\hfill\okyten{\dauSinhVien}{Sinh viên thực hiện}

\end{acknowledgments}

% ---------------------------------------------------------------------
%  TÓM TẮT
% ---------------------------------------------------------------------
\begin{abstract}

Tìm phòng trọ phù hợp là nhu cầu thiết yếu của sinh viên và người lao động
tại Đà Nẵng, nhưng thông tin cho thuê hiện nay nằm rải rác trên nhiều kênh,
thiếu công cụ lọc theo đúng nhu cầu cá nhân về ngân sách, khu vực và tiện
ích. Đồ án xây dựng một hệ thống hỗ trợ tìm phòng trọ tại Đà Nẵng, trong đó
thành phần cốt lõi là một hệ gợi ý theo hướng dựa trên nội dung
\en{(content-based filtering)}, kết hợp thêm tín hiệu từ những người dùng có
nhu cầu tương tự để tăng độ đa dạng của gợi ý. Hệ thống trích xuất đặc trưng
của từng phòng trọ gồm giá thuê, khu vực, diện tích, tiện ích và mô tả văn
bản, sau đó tính độ tương đồng có trọng số với hồ sơ nhu cầu người dùng để
xếp hạng kết quả. Dữ liệu thử nghiệm gồm 2.500 tin đăng phòng trọ thu thập
tại sáu quận trung tâm của Đà Nẵng, cùng lịch sử tương tác của 200 người
dùng thử nghiệm dùng để đánh giá ngoại tuyến. Hệ thống đạt Precision@5 bằng
0,76 và NDCG@5 bằng 0,81, cao hơn mô hình chỉ dùng nội dung thuần tuý lần
lượt 7 và 6 điểm phần trăm. Khảo sát nhanh trên 30 sinh viên dùng thử cho
mức độ hài lòng trung bình 4,1 trên thang 5. Toàn bộ hệ thống được triển
khai thành ứng dụng web cho phép người dùng nhập tiêu chí và nhận danh sách
phòng trọ gợi ý kèm bản đồ vị trí. Kết quả cho thấy cách tiếp cận dựa trên
nội dung là lựa chọn khả thi khi dữ liệu tương tác người dùng còn hạn chế,
và có thể mở rộng thêm thành phần cộng tác khi hệ thống có nhiều người dùng
hơn.

\vspace{0.8em}
\noindent\textbf{Từ khoá:} hệ gợi ý, lọc dựa trên nội dung, độ tương đồng
cô-sin, phòng trọ, Đà Nẵng.

\end{abstract}


\tableofcontents
\listoffigures
\listoftables
% Bỏ dòng dưới nếu đồ án của em không có giải thuật nào
\listofalgorithms

% ---------------------------------------------------------------------
%  THÂN BÁO CÁO: sáu chương theo quy cách của học phần
% ---------------------------------------------------------------------
\mainmatter
\dauPageStyle

% ---------------------------------------------------------------------
%  CHƯƠNG 1 — MỞ ĐẦU
% ---------------------------------------------------------------------
\chapter{Mở đầu}
\label{ch:modau}

\begin{muctieu}
Chương mở đầu trình bày lý do chọn đề tài, mục tiêu cần đạt, phạm vi thực
hiện và bố cục báo cáo, giúp người đọc nắm được bức tranh tổng quan về đồ án
xây dựng hệ thống hỗ trợ tìm phòng trọ tại Đà Nẵng.
\end{muctieu}

\section{Lý do chọn đề tài}

Đà Nẵng là một trong những thành phố có số lượng sinh viên và lao động
ngoại tỉnh nhập cư đông, kéo theo nhu cầu thuê phòng trọ rất lớn quanh các
khu vực trường đại học, khu công nghiệp và trung tâm thành phố. Tuy nhiên,
thông tin cho thuê phòng trọ hiện nằm rải rác trên nhiều kênh khác nhau như
mạng xã hội, các nhóm chia sẻ và biển quảng cáo dán trước nhà, thiếu tính
thống nhất và khó tra cứu theo đúng nhu cầu cá nhân. Người tìm trọ thường
phải xem qua hàng chục tin đăng, so sánh thủ công giữa giá thuê, vị trí và
tiện ích, mất nhiều thời gian mà vẫn có thể bỏ sót những lựa chọn phù hợp.
Bài toán gợi ý \en{(recommendation)} trong trí tuệ nhân tạo cho phép tự
động xếp hạng và đề xuất những phòng trọ phù hợp nhất dựa trên tiêu chí của
từng người dùng, thay vì buộc người dùng phải tự lọc qua toàn bộ danh sách.
Đề tài này được chọn vì vừa giải quyết một nhu cầu thực tế gần gũi với đời
sống sinh viên, vừa là cơ hội để em vận dụng kiến thức về hệ gợi ý, xử lý
dữ liệu và xây dựng ứng dụng web đã học trong học phần.

\section{Mục tiêu của đồ án}

Đồ án hướng tới các mục tiêu sau:
\begin{itemize}[leftmargin=1.6em, itemsep=2pt]
  \item Xây dựng được bộ dữ liệu phòng trọ tại khu vực Đà Nẵng với đầy đủ
        thông tin về giá thuê, vị trí, diện tích và tiện ích.
  \item Cài đặt được một mô hình gợi ý phòng trọ dựa trên nội dung, có khả
        năng xếp hạng phòng trọ theo mức độ phù hợp với tiêu chí người
        dùng nhập vào.
  \item Đạt Precision@5 từ 0,7 trở lên khi đánh giá ngoại tuyến trên tập
        tương tác người dùng thử nghiệm.
  \item Phát triển giao diện web cho phép người dùng nhập tiêu chí tìm
        phòng và xem danh sách gợi ý kèm bản đồ vị trí.
  \item So sánh mô hình gợi ý dựa trên nội dung thuần tuý với phiên bản có
        kết hợp tín hiệu từ người dùng tương tự, để làm rõ mức độ cải thiện.
\end{itemize}

\section{Phạm vi và đối tượng}

Đồ án tập trung vào việc gợi ý phòng trọ, nhà trọ cho thuê dài hạn dành cho
sinh viên và người đi làm tại khu vực nội thành Đà Nẵng, cụ thể là sáu
quận trung tâm gồm Hải Châu, Thanh Khê, Sơn Trà, Ngũ Hành Sơn, Cẩm Lệ và
Liên Chiểu. Dữ liệu phòng trọ được thu thập thông qua khảo sát trực tiếp và
từ những tin đăng công khai được chủ trọ đồng ý cung cấp. Đồ án không xử lý
loại hình căn hộ chung cư cao cấp hay khách sạn, không thực hiện chức năng
đặt cọc hay thanh toán trực tuyến, và chưa triển khai ở quy mô sản xuất
phục vụ số lượng người dùng lớn.

\begin{luuy}
Lỗi phổ biến nhất của đồ án là chọn phạm vi quá rộng rồi bỏ dở. Một bài toán
hẹp làm tới nơi tới chốn luôn được chấm cao hơn một ý tưởng lớn còn dang dở.
\end{luuy}

\section{Bố cục báo cáo}

Phần còn lại của báo cáo được tổ chức như sau. Chương~\ref{ch:coso} trình
bày cơ sở lý thuyết và công nghệ sử dụng. Chương~\ref{ch:thietke} phân tích
bài toán và thiết kế giải pháp. Chương~\ref{ch:caidat} mô tả quá trình cài
đặt và thử nghiệm. Chương~\ref{ch:ketqua} trình bày kết quả và đánh giá.
Chương~\ref{ch:ketluan} tổng kết và nêu hướng phát triển.

% ---------------------------------------------------------------------
%  CHƯƠNG 2 — CƠ SỞ LÝ THUYẾT VÀ CÔNG NGHỆ
% ---------------------------------------------------------------------
\chapter{Cơ sở lý thuyết và công nghệ}
\label{ch:coso}

\begin{muctieu}
Chương này trình bày lý thuyết về hệ gợi ý và cách tính độ tương đồng dùng
trong đồ án, cùng danh mục công cụ và thư viện để xây dựng hệ thống.
\end{muctieu}

\section{Kiến thức nền tảng}

\begin{dinhnghia}[Hệ gợi ý]
Hệ gợi ý \en{(recommender system)} là hệ thống dự đoán mức độ phù hợp giữa
một người dùng $u$ và một đối tượng $i$ trong tập đối tượng $\mathcal{I}$,
từ đó xếp hạng và đề xuất những đối tượng có điểm phù hợp cao nhất cho
người dùng.
\end{dinhnghia}

Hai hướng tiếp cận phổ biến nhất là lọc dựa trên nội dung
\en{(content-based filtering)}, dựa vào đặc trưng của chính đối tượng và hồ
sơ nhu cầu của người dùng, và lọc cộng tác \en{(collaborative filtering)},
dựa vào hành vi của những người dùng có sở thích tương tự. Với bài toán tìm
phòng trọ, số lượt tương tác của mỗi người dùng thường rất ít vì nhu cầu
thuê trọ không lặp lại thường xuyên, nên lọc cộng tác thuần tuý dễ gặp vấn
đề khởi động lạnh \en{(cold start)}. Đồ án chọn lọc dựa trên nội dung làm
thành phần chính, vì mỗi phòng trọ luôn có đầy đủ đặc trưng mô tả ngay khi
được đăng, không phụ thuộc vào lịch sử tương tác.

Mỗi phòng trọ $i$ được biểu diễn bằng một vec-tơ đặc trưng $\mathbf{x}_i$
gồm các thành phần đã chuẩn hoá: giá thuê, khu vực, diện tích, danh sách
tiện ích và mô tả văn bản vec-tơ hoá theo TF-IDF. Nhu cầu của người dùng
cũng được biểu diễn bằng một vec-tơ $\mathbf{q}$ cùng cấu trúc, ghép từ
ngân sách mong muốn, khu vực ưu tiên, diện tích tối thiểu và các tiện ích
bắt buộc. Độ phù hợp giữa nhu cầu và một phòng trọ được tính bằng tổng có
trọng số của độ tương đồng cô-sin trên từng nhóm đặc trưng
\begin{equation}
\mathrm{score}(\mathbf{q}, \mathbf{x}_i) =
  w_1 \cdot \mathrm{sim}_{\text{giá}} +
  w_2 \cdot \mathrm{sim}_{\text{khu vực}} +
  w_3 \cdot \mathrm{sim}_{\text{tiện ích}} +
  w_4 \cdot \mathrm{sim}_{\text{mô tả}},
\label{eq:score}
\end{equation}
trong đó mỗi $\mathrm{sim}(\cdot)$ là độ tương đồng cô-sin
\begin{equation}
\mathrm{sim}(\mathbf{a}, \mathbf{b}) =
  \frac{\mathbf{a} \cdot \mathbf{b}}{\lVert \mathbf{a} \rVert \, \lVert \mathbf{b} \rVert}
\label{eq:cosine}
\end{equation}
tính trên nhóm thành phần tương ứng, và $w_1, \dots, w_4$ là trọng số không
âm có tổng bằng một. Các phòng trọ được xếp hạng giảm dần theo điểm số ở
Công thức~\eqref{eq:score} và trả về $k$ phòng có điểm cao nhất.

\begin{ghinho}
Việc chuẩn hoá từng nhóm đặc trưng về cùng thang giá trị trước khi tính độ
tương đồng cô-sin là bắt buộc, nếu không thành phần có giá trị tuyệt đối
lớn hơn, ví dụ giá thuê tính bằng đồng, sẽ lấn át hoàn toàn các thành phần
còn lại trong Công thức~\eqref{eq:score}.
\end{ghinho}

\section{Công cụ và thư viện}

\begin{table}[H]
\centering
\caption{Công cụ và thư viện sử dụng trong đồ án.}
\label{tab:congcu}
\begin{tabular*}{\textwidth}{@{\extracolsep{\fill}}>{\bfseries\color{daunavy}}l >{\color{daumaroon}}l l}
\toprule
\rowcolor{gray100}\bfseries\color{daunavy}Thành phần & \bfseries\color{daumaroon}Phiên bản & \bfseries\color{daunavy}Vai trò \\
\midrule
Python       & 3.11 & Ngôn ngữ lập trình chính \\
pandas       & 2.2  & Đọc và xử lý dữ liệu phòng trọ dạng bảng \\
scikit-learn & 1.4  & Vec-tơ hoá TF-IDF và tính độ tương đồng cô-sin \\
Flask        & 3.0  & Xây dựng giao diện web và API gợi ý \\
SQLite       & 3.x  & Lưu trữ dữ liệu phòng trọ và người dùng \\
Leaflet.js   & 1.9  & Hiển thị bản đồ vị trí phòng trọ trên nền OpenStreetMap \\
\bottomrule
\end{tabular*}
\end{table}

\section{Các công trình liên quan}

Các hệ thống gợi ý bất động sản và phòng cho thuê thường đi theo hai hướng
chính. Hướng thứ nhất dùng lọc dựa trên nội dung với đặc trưng thủ công như
giá, vị trí và diện tích, có ưu điểm dễ triển khai và không cần dữ liệu
tương tác lớn~\cite{RussellNorvig2021}. Hướng thứ hai kết hợp thêm lọc cộng
tác để tận dụng hành vi của những người dùng có sở thích tương tự, giúp
tăng tính đa dạng của kết quả gợi ý nhưng đòi hỏi lượng dữ liệu tương tác đủ
lớn để tránh khởi động lạnh~\cite{HoangKiem}. Đồ án này chọn hướng thứ nhất
làm nền tảng, đồng thời bổ sung một thành phần cộng tác đơn giản dựa trên
những phòng trọ mà người dùng có hồ sơ nhu cầu tương tự đã lưu lại, nhằm
giảm bớt hạn chế của lọc dựa trên nội dung thuần tuý là kết quả gợi ý có xu
hướng quá giống với tiêu chí đã nhập, thiếu bất ngờ hữu ích.

% ---------------------------------------------------------------------
%  CHƯƠNG 3 — PHÂN TÍCH VÀ THIẾT KẾ
% ---------------------------------------------------------------------
\chapter{Phân tích và thiết kế}
\label{ch:thietke}

\begin{muctieu}
Chương này phân tích yêu cầu bài toán, đề xuất giải pháp, thiết kế kiến trúc
hệ thống, thiết kế dữ liệu và thiết kế mô hình gợi ý phòng trọ.
\end{muctieu}

\section{Phân tích bài toán và yêu cầu}

Hệ thống nhận đầu vào là tiêu chí tìm phòng của người dùng, gồm ngân sách
tối đa, khu vực mong muốn, diện tích tối thiểu và danh sách tiện ích bắt
buộc, và trả về danh sách phòng trọ được xếp hạng theo mức độ phù hợp kèm
điểm số. Yêu cầu chức năng gồm nhập tiêu chí tìm kiếm, tính điểm và xếp
hạng phòng trọ, hiển thị kết quả kèm bản đồ vị trí, và cho phép người dùng
lưu lại phòng trọ quan tâm để cải thiện gợi ý ở những lần tìm kiếm sau. Yêu
cầu phi chức năng gồm thời gian phản hồi dưới hai giây cho mỗi lượt tìm
kiếm trên tập dữ liệu hiện có, và giao diện dùng được trên trình duyệt phổ
thông cả ở máy tính lẫn điện thoại.

\section{Thiết kế kiến trúc hệ thống}

Kiến trúc hệ thống được trình bày ở Hình~\ref{fig:kientruc}. Dữ liệu đi
tuần tự qua bốn khối, mỗi khối là một mô-đun độc lập trong mã nguồn.

\begin{figure}[H]
\centering
\resizebox{\textwidth}{!}{%
\begin{tikzpicture}[node distance=12mm and 14mm]
  \node[objbox] (in) {Tiêu chí\\ tìm phòng};
  \node[clsbox, right=of in]  (pre) {Chuẩn hoá\\ và vec-tơ hoá};
  \node[clsbox, right=of pre] (sim) {Tính độ\\ tương đồng};
  \node[clsbox, right=of sim] (rank) {Xếp hạng\\ theo điểm số};
  \node[objbox, right=of rank] (out) {Danh sách\\ phòng gợi ý};
  \draw[flowarr] (in)   -- (pre);
  \draw[flowarr] (pre)  -- (sim);
  \draw[flowarr] (sim)  -- (rank);
  \draw[flowarr] (rank) -- (out);
\end{tikzpicture}}
\caption{Sơ đồ kiến trúc tổng thể của hệ thống.}
\label{fig:kientruc}
\end{figure}

Khối tiêu chí tìm phòng nhận đầu vào từ giao diện web. Khối chuẩn hoá và
vec-tơ hoá đưa tiêu chí người dùng và hồ sơ từng phòng trọ về cùng thang
giá trị. Khối tính độ tương đồng áp dụng Công thức~\eqref{eq:score} trên
toàn bộ phòng trọ còn trống trong cơ sở dữ liệu. Khối xếp hạng chọn ra
$k$ phòng có điểm cao nhất, đồng thời trộn thêm một số phòng trọ được người
dùng có hồ sơ tương tự đã lưu lại, trước khi trả kết quả cho giao diện.

\section{Thiết kế dữ liệu}

Bộ dữ liệu gồm 2.500 tin đăng phòng trọ tại sáu quận trung tâm Đà Nẵng,
được chia thành 2.000 phòng dùng để xây dựng hồ sơ đặc trưng phục vụ gợi ý
và 500 phòng giữ lại làm tập kiểm thử ngoại tuyến, cùng lịch sử tương tác
đã lưu hoặc đã liên hệ của 200 người dùng thử nghiệm. Cấu trúc mỗi bản ghi
phòng trọ được mô tả ở Bảng~\ref{tab:dulieu}.

\begin{table}[H]
\centering
\caption{Cấu trúc bản ghi dữ liệu phòng trọ.}
\label{tab:dulieu}
\begin{tabular*}{\textwidth}{@{\extracolsep{\fill}}>{\bfseries\color{daunavy}}l >{\color{daumaroon}}l l}
\toprule
\rowcolor{gray100}\bfseries\color{daunavy}Trường & \bfseries\color{daumaroon}Kiểu & \bfseries\color{daunavy}Mô tả \\
\midrule
id          & integer & Khoá định danh phòng trọ \\
gia\_thue   & integer & Giá thuê hằng tháng, đơn vị nghìn đồng \\
quan        & text    & Quận nơi phòng trọ toạ lạc \\
dien\_tich  & float   & Diện tích phòng, đơn vị mét vuông \\
tien\_ich   & text    & Danh sách tiện ích, phân tách bằng dấu phẩy \\
mo\_ta      & text    & Mô tả chi tiết do chủ trọ cung cấp \\
vi\_do, kinh\_do & float & Toạ độ dùng để hiển thị trên bản đồ \\
\bottomrule
\end{tabular*}
\end{table}

\section{Thiết kế mô hình gợi ý}

Mô hình được chọn là lọc dựa trên nội dung có trọng số theo
Công thức~\eqref{eq:score}, với bốn nhóm đặc trưng: giá thuê và diện tích
được chuẩn hoá theo Min-Max, khu vực được mã hoá one-hot theo quận, tiện
ích được mã hoá one-hot theo từng loại tiện ích, và mô tả văn bản được
vec-tơ hoá bằng TF-IDF với n-gram bậc một. Trọng số $w_1, \dots, w_4$ trong
Công thức~\eqref{eq:score} được chọn bằng tìm kiếm lưới trên tập kiểm thử
ngoại tuyến, tối ưu theo Precision@5, không đụng tới 500 phòng trọ giữ lại
làm tập đánh giá cuối cùng trong quá trình chọn trọng số. Thành phần cộng
tác được cài đặt đơn giản bằng cách tìm năm người dùng có vec-tơ tiêu chí
gần nhất với người dùng hiện tại, sau đó cộng thêm một phần điểm cho những
phòng trọ đã được các người dùng đó lưu lại.

\begin{luuy}
Không được dùng tập kiểm thử cuối cùng để chọn trọng số hay siêu tham số.
Làm vậy thì các độ đo báo cáo ở Chương~\ref{ch:ketqua} sẽ cao hơn thực tế,
và em sẽ không giải thích được khi hội đồng hỏi tới.
\end{luuy}

% ---------------------------------------------------------------------
%  CHƯƠNG 4 — CÀI ĐẶT VÀ THỬ NGHIỆM
% ---------------------------------------------------------------------
\chapter{Cài đặt và thử nghiệm}
\label{ch:caidat}

\begin{muctieu}
Chương này mô tả môi trường triển khai, các thành phần chính kèm mã nguồn
minh hoạ, và giao diện demo của hệ thống gợi ý phòng trọ.
\end{muctieu}

\section{Môi trường triển khai}

Hệ thống được phát triển và thử nghiệm trên máy tính cá nhân với CPU Intel
Core i5, RAM 8~GB, hệ điều hành Windows 11 và Python 3.11 cùng các thư viện
đã nêu ở Bảng~\ref{tab:congcu}. Mã nguồn được quản lý bằng \code{git} và đặt
tại kho \url{https://github.com/tencuaem/tim-phong-tro-da-nang}.

\section{Các thành phần chính}

Mã nguồn~\ref{lst:featurize} minh hoạ bước chuẩn hoá và vec-tơ hoá đặc
trưng phòng trọ.

\begin{lstlisting}[language=Python, caption={Chuẩn hoá và vec-tơ hoá đặc trưng phòng trọ.}, label={lst:featurize}]
from sklearn.preprocessing import MinMaxScaler, OneHotEncoder
from sklearn.feature_extraction.text import TfidfVectorizer

# Normalize numeric features to the same [0, 1] scale
price_area_scaler = MinMaxScaler()
price_area = price_area_scaler.fit_transform(df[["gia_thue", "dien_tich"]])

# One-hot encode categorical district and amenity list
district_encoder = OneHotEncoder(sparse_output=False)
district_onehot = district_encoder.fit_transform(df[["quan"]])

# TF-IDF for the free-text description field
tfidf = TfidfVectorizer(max_features=500)
desc_vectors = tfidf.fit_transform(df["mo_ta"])
\end{lstlisting}

Mã nguồn~\ref{lst:score} minh hoạ hàm tính điểm gợi ý theo
Công thức~\eqref{eq:score} và xếp hạng phòng trọ.

\begin{lstlisting}[language=Python, caption={Tính điểm gợi ý và xếp hạng phòng trọ.}, label={lst:score}]
from sklearn.metrics.pairwise import cosine_similarity

# Weighted score across the four feature groups (see Equation 2.1)
def score_rooms(query_vec, room_vecs, weights):
    sim_price = cosine_similarity(query_vec["price_area"], room_vecs["price_area"])
    sim_area = cosine_similarity(query_vec["district"], room_vecs["district"])
    sim_amenity = cosine_similarity(query_vec["amenity"], room_vecs["amenity"])
    sim_desc = cosine_similarity(query_vec["desc"], room_vecs["desc"])

    w1, w2, w3, w4 = weights
    return w1 * sim_price + w2 * sim_area + w3 * sim_amenity + w4 * sim_desc

def top_k_rooms(scores, room_ids, k=5):
    # Return the k highest-scoring room ids, descending order
    ranked = sorted(zip(room_ids, scores), key=lambda x: x[1], reverse=True)
    return ranked[:k]
\end{lstlisting}

\section{Giao diện và demo}

\begin{figure}[H]
\centering
% Bỏ dấu % ở dòng dưới và thay ten-anh.png bằng ảnh thật của em
% \includegraphics[width=0.8\textwidth]{ten-anh.png}
\fbox{\begin{minipage}[c][5cm][c]{0.8\textwidth}\centering
  \itshape\color{gray500}[Chèn ảnh chụp màn hình giao diện tìm phòng trọ tại đây]
\end{minipage}}
\caption{Giao diện web cho phép nhập tiêu chí và xem danh sách phòng trọ
gợi ý kèm bản đồ vị trí.}
\label{fig:demo}
\end{figure}

\section{Kiểm thử}

Hệ thống được kiểm thử với ba nhóm ca. Nhóm một là tiêu chí thông thường
với ngân sách và khu vực rõ ràng. Nhóm hai là tiêu chí hẹp, ví dụ ngân sách
rất thấp kết hợp yêu cầu nhiều tiện ích, để kiểm tra hệ thống xử lý thế nào
khi không có phòng trọ nào thoả mãn đầy đủ. Nhóm ba là dữ liệu biên gồm
không nhập tiêu chí nào, nhập ngân sách âm và nhập khu vực không tồn tại
trong dữ liệu. Ở nhóm ba, hệ thống ban đầu báo lỗi khi ngân sách âm; lỗi
được khắc phục bằng cách kiểm tra và chặn giá trị không hợp lệ trước khi
đưa vào Công thức~\eqref{eq:score}.

% ---------------------------------------------------------------------
%  CHƯƠNG 5 — KẾT QUẢ VÀ ĐÁNH GIÁ
% ---------------------------------------------------------------------
\chapter{Kết quả và đánh giá}
\label{ch:ketqua}

\begin{muctieu}
Chương này trình bày kết quả đạt được, các độ đo đánh giá định lượng, phần
bàn luận và những hạn chế còn tồn tại của hệ thống gợi ý phòng trọ.
\end{muctieu}

\section{Kết quả đạt được}

Hệ thống đã hoàn thành các mục tiêu đề ra ở Chương~\ref{ch:modau}. Bộ dữ
liệu 2.500 phòng trọ tại sáu quận trung tâm Đà Nẵng đã được xây dựng và làm
sạch, mô hình gợi ý dựa trên nội dung đạt Precision@5 bằng 0,76, vượt
ngưỡng mục tiêu 0,7, và giao diện web cho phép người dùng nhập tiêu chí và
nhận danh sách phòng trọ gợi ý kèm bản đồ trong thời gian thực.

\section{Độ đo đánh giá}

Vì bài toán là xếp hạng và gợi ý chứ không phải phân lớp, đồ án dùng các độ
đo phổ biến cho hệ gợi ý gồm Precision@5, Recall@5 và NDCG@5, tính trên
lịch sử tương tác đã lưu hoặc đã liên hệ của 200 người dùng thử nghiệm, đối
chiếu với 500 phòng trọ ở tập kiểm thử ngoại tuyến. Bảng~\ref{tab:ketqua}
so sánh mô hình chỉ dùng lọc dựa trên nội dung với phiên bản kết hợp thêm
thành phần cộng tác đơn giản.

\begin{table}[H]
\centering
\caption{So sánh hiệu năng các phiên bản mô hình gợi ý.}
\label{tab:ketqua}
\begin{tabular*}{\textwidth}{@{\extracolsep{\fill}}>{\bfseries\color{daunavy}}l ccc}
\toprule
\rowcolor{gray100}\bfseries\color{daunavy}Mô hình & \bfseries\color{daumaroon}Precision@5 & \bfseries\color{daumaroon}Recall@5 & \bfseries\color{daumaroon}NDCG@5 \\
\midrule
Chỉ lọc dựa trên nội dung          & 0,69 & 0,60 & 0,75 \\
Kết hợp thành phần cộng tác        & 0,76 & 0,68 & 0,81 \\
\bottomrule
\end{tabular*}
\end{table}

Ngoài đánh giá ngoại tuyến, đồ án tổ chức khảo sát nhanh trên 30 sinh viên
dùng thử hệ thống, chấm mức độ hài lòng theo thang điểm 5. Kết quả trung
bình đạt 4,1 điểm, trong đó tiêu chí được đánh giá cao nhất là tốc độ phản
hồi và tiêu chí thấp nhất là độ đa dạng của kết quả gợi ý.

\section{Bàn luận}

Phiên bản kết hợp thành phần cộng tác nhỉnh hơn phiên bản chỉ dùng nội dung
khoảng bảy điểm phần trăm ở Precision@5. Mức cải thiện này đến từ việc một
số phòng trọ tốt nhưng có mô tả sơ sài, dẫn đến độ tương đồng văn bản thấp,
vẫn được gợi ý nhờ được người dùng có nhu cầu tương tự lưu lại trước đó.
Phân tích các trường hợp gợi ý chưa tốt cho thấy phần lớn tập trung ở những
tiêu chí có khu vực mong muốn ít phòng trọ, khi đó hệ thống buộc phải nới
lỏng điều kiện khu vực để trả đủ năm kết quả, làm giảm độ phù hợp trung
bình của danh sách gợi ý.

\section{Hạn chế}

\begin{luuy}
Hệ thống còn gợi ý kém đa dạng khi tiêu chí người dùng quá chặt, do bản chất
lọc dựa trên nội dung có xu hướng trả về những phòng trọ giống nhau. Thành
phần cộng tác chỉ dựa trên năm người dùng gần nhất nên chưa phát huy hết
tác dụng khi số lượng người dùng còn ít. Dữ liệu phòng trọ cũng chưa được
cập nhật tự động, một số tin đăng có thể đã hết phòng nhưng vẫn còn trong
hệ thống.
\end{luuy}

% ---------------------------------------------------------------------
%  CHƯƠNG 6 — KẾT LUẬN VÀ HƯỚNG PHÁT TRIỂN
% ---------------------------------------------------------------------
\chapter{Kết luận và hướng phát triển}
\label{ch:ketluan}

\begin{muctieu}
Chương này tổng kết những gì đồ án đã làm được, rút ra bài học và đề xuất
hướng phát triển tiếp theo cho hệ thống gợi ý phòng trọ.
\end{muctieu}

\section{Kết luận}

Đồ án đã xây dựng thành công một hệ thống hỗ trợ tìm phòng trọ tại Đà Nẵng,
đi trọn từ khâu thu thập và chuẩn hoá dữ liệu, thiết kế mô hình gợi ý dựa
trên nội dung kết hợp thành phần cộng tác đơn giản, đến triển khai giao
diện web có bản đồ vị trí. Qua đó em đã vận dụng được kiến thức về hệ gợi
ý, xử lý dữ liệu và xây dựng ứng dụng web vào một bài toán gắn liền với đời
sống sinh viên, đồng thời làm quen với quy trình phát triển phần mềm có
quản lý phiên bản và đánh giá định lượng bằng các độ đo chuyên biệt cho
bài toán xếp hạng.

\section{Bài học rút ra}

Bài học lớn nhất là việc thiết kế đặc trưng và chọn trọng số giữa các nhóm
đặc trưng ảnh hưởng đến chất lượng gợi ý nhiều hơn việc chọn công thức tính
độ tương đồng. Phần lớn thời gian của đồ án dành cho thu thập, làm sạch dữ
liệu phòng trọ và thử nghiệm các cách chuẩn hoá khác nhau, trong khi việc
đổi công thức tính điểm chỉ đem lại thay đổi nhỏ. Em cũng nhận ra hạn chế
của cách tiếp cận thuần dựa trên nội dung là dễ khiến kết quả gợi ý thiếu
đa dạng, và điều này chỉ được khắc phục một phần bằng thành phần cộng tác
đơn giản. Nếu làm lại, em sẽ dành thời gian đầu để thu thập dữ liệu tương
tác người dùng nhiều hơn, thay vì tập trung phần lớn công sức vào riêng dữ
liệu phòng trọ.

\section{Hướng phát triển}

Một số hướng phát triển tiếp theo:
\begin{itemize}[leftmargin=1.6em, itemsep=2pt]
  \item Mở rộng thành phần cộng tác bằng các phương pháp phân rã ma trận,
        nhằm khắc phục hạn chế về độ đa dạng đã nêu ở Chương~\ref{ch:ketqua}.
  \item Bổ sung cơ chế cập nhật tự động trạng thái còn trống hay đã cho
        thuê của từng phòng trọ, tránh gợi ý những tin đăng đã hết hạn.
  \item Mở rộng dữ liệu ra toàn bộ các quận huyện của Đà Nẵng và thu thập
        thêm lịch sử tương tác thực tế để cải thiện thành phần cộng tác.
  \item Triển khai hệ thống lên nền tảng đám mây và bổ sung ứng dụng di
        động để phục vụ nhiều người dùng đồng thời.
\end{itemize}


% Chương tra cứu quy cách trình bày. XOÁ DÒNG DƯỚI khi nộp bản chính
% thức, vì đây chỉ là sổ tra cứu chứ không phải một chương của báo cáo.
% ---------------------------------------------------------------------
%  CHƯƠNG 2 — QUY CÁCH TRÌNH BÀY
%  Chương này đóng vai trò sổ tra cứu: mỗi mục là một mẫu em copy rồi sửa.
%  Khi viết đồ án thật, em thay cả chương này bằng nội dung của mình.
% ---------------------------------------------------------------------
\chapter{Quy cách trình bày trong báo cáo}
\label{ch:soanthao}

\begin{muctieu}
Chương này gom sẵn mẫu cho mọi thành phần em cần dùng: hình, bảng, công
thức, mã nguồn, giải thuật, trích dẫn, chú thích và các hộp nhấn mạnh. Mỗi
mục có đoạn mã mẫu, em copy rồi thay nội dung.
\end{muctieu}

\begin{luuy}
Khi bắt đầu viết đồ án thật, hãy xoá cả chương này. Nó chỉ là sổ tra cứu,
không phải một chương của báo cáo.
\end{luuy}

\section{Chèn hình ảnh}

Đặt file ảnh vào thư mục \code{figs/} rồi dùng mẫu dưới đây. Mọi hình đều
phải có chú thích, có nhãn, và phải được nhắc tới trong thân bài như
Hình~\ref{fig:vidu}.

\begin{figure}[H]
\centering
\includegraphics[width=0.72\textwidth]{dau-campus.png}
\caption{Một ví dụ về chèn hình ảnh từ file.}
\label{fig:vidu}
\end{figure}

Nếu hình là sơ đồ do em tự vẽ thì nên vẽ thẳng bằng TikZ như
Hình~\ref{fig:luongxuly}, vì sơ đồ vẽ bằng TikZ luôn nét khi phóng to và
dùng đúng bộ màu của trường.

\begin{figure}[H]
\centering
% \resizebox ép sơ đồ vừa đúng bề ngang trang, tránh tràn lề
\resizebox{\textwidth}{!}{%
\begin{tikzpicture}[node distance=8mm and 12mm]
  \node[io] (in) {Ảnh đầu vào};
  \node[process, right=of in] (pre) {Tiền xử lý};
  \node[decision, right=of pre] (chk) {Đủ chất lượng?};
  \node[process, right=of chk] (mod) {Mô hình CNN};
  \node[startstop, below=of chk] (loai) {Loại bỏ};
  \draw[flowarr] (in)  -- (pre);
  \draw[flowarr] (pre) -- (chk);
  \draw[flowarr] (chk) -- node[above, font=\scriptsize] {có} (mod);
  \draw[flowarr] (chk) -- node[right, font=\scriptsize] {không} (loai);
\end{tikzpicture}}
\caption{Một ví dụ sơ đồ luồng xử lý vẽ bằng TikZ.}
\label{fig:luongxuly}
\end{figure}

\section{Chèn bảng}

Bảng dùng chú thích đặt phía trên, đúng quy cách trình bày khoa học. Dùng
\code{booktabs} với ba loại đường kẻ \code{toprule}, \code{midrule},
\code{bottomrule}, tránh kẻ dọc cho mọi ô.

\begin{table}[H]
\centering
\caption{Một ví dụ về bảng so sánh kết quả.}
\label{tab:sosanh}
\begin{tabular*}{\textwidth}{@{\extracolsep{\fill}}>{\bfseries\color{daunavy}}l cccc}
\toprule
\rowcolor{gray100}\bfseries\color{daunavy}Mô hình & \bfseries\color{daumaroon}Accuracy & \bfseries\color{daumaroon}Precision & \bfseries\color{daumaroon}Recall & \bfseries\color{daumaroon}F1 \\
\midrule
Huấn luyện từ đầu   & 0,848 & 0,841 & 0,836 & 0,838 \\
Học chuyển giao     & 0,912 & 0,907 & 0,903 & 0,905 \\
\bottomrule
\end{tabular*}
\end{table}

\begin{luuy}
Bảng rộng quá lề là lỗi hay gặp nhất. Nếu bảng tràn, hãy giảm cỡ chữ bằng
\code{small} hoặc \code{footnotesize}, hoặc gộp bớt cột. Đừng để LaTeX tự
tràn rồi bỏ qua.
\end{luuy}

\section{Chèn công thức toán học}

Công thức đặt giữa dòng và có đánh số thì dùng môi trường \code{equation},
như Công thức~\eqref{eq:softmax}. Công thức ngắn chèn trong câu thì viết
giữa hai dấu đô la, ví dụ $y = f(x)$.

\begin{equation}
p(y = k \mid \mathbf{x}) = \frac{\exp(z_k)}{\sum_{j=1}^{K} \exp(z_j)},
\label{eq:softmax}
\end{equation}

trong đó $z_k$ là điểm số thô mà mạng gán cho lớp $k$, và $K$ là số lớp. Mọi
ký hiệu xuất hiện trong công thức đều phải được giải thích ngay sau đó, như
câu này đang làm với \eqref{eq:softmax}.

\begin{ghinho}
Công thức là một phần của câu văn, không phải một khối trôi nổi. Trước công
thức phải có câu dẫn vào, sau công thức phải có câu giải thích ký hiệu hoặc
ý nghĩa. Công thức đứng một mình không ai đọc.
\end{ghinho}

\section{Chèn mã nguồn}

Mẫu này dùng gói \code{listings} thay cho \code{minted} của bản gốc, vì
\code{listings} không cần bật \code{-shell-escape} nên build được ngay trên
Overleaf lẫn máy cá nhân.

Cách thứ nhất là gõ trực tiếp, kết quả như Mã nguồn~\ref{lst:tructiep}.

\begin{lstlisting}[language=Python, caption={Chèn mã nguồn trực tiếp vào báo cáo.}, label={lst:tructiep}]
import torch.nn as nn

# A small convolutional block: conv, batch norm, activation
def conv_block(in_ch, out_ch):
    return nn.Sequential(
        nn.Conv2d(in_ch, out_ch, kernel_size=3, padding=1),
        nn.BatchNorm2d(out_ch),
        nn.ReLU(inplace=True),
    )
\end{lstlisting}

Cách thứ hai là nạp mã từ một file có sẵn trong đồ án. Cách này tiện vì báo
cáo luôn khớp với mã thật, sửa mã là báo cáo tự cập nhật theo. Mã
nguồn~\ref{lst:tufile} được nạp từ file \code{code/vi-du-huan-luyen.py}.

\lstinputlisting[language=Python, caption={Chèn mã nguồn từ file có sẵn.}, label={lst:tufile}]{code/vi-du-huan-luyen.py}

\begin{luuy}
Quy ước của khoa: comment trong code viết bằng tiếng Anh, còn văn xuôi của
báo cáo viết bằng tiếng Việt. Không đặt tiếng Việt bên trong khối mã nguồn,
vì phông mã nguồn không có đủ dấu tiếng Việt.
\end{luuy}

Chỉ đưa vào thân bài những đoạn mã cốt lõi. Mã dài để ở phụ lục hoặc chỉ dẫn
tới kho \code{git} của đồ án.

\section{Biểu diễn giải thuật}

Giải thuật trình bày bằng mã giả, không phải mã của một ngôn ngữ cụ thể,
để người đọc nắm được ý tưởng mà không cần biết ngôn ngữ đó.
Giải thuật~\ref{alg:nguyento} là một ví dụ.

\begin{algorithm}[H]
\caption{Kiểm tra một số tự nhiên có phải số nguyên tố hay không}
\label{alg:nguyento}
\hspace*{\algorithmicindent}\textbf{Đầu vào:} số tự nhiên $n$\\
\hspace*{\algorithmicindent}\textbf{Đầu ra:} \textsc{True} nếu $n$ là số nguyên tố, ngược lại \textsc{False}
\begin{algorithmic}[1]
\Function{LaSoNguyenTo}{$n$}
  \If{$n < 2$}
    \State \Return \textsc{False}
  \EndIf
  \For{$i \gets 2$ \textbf{to} $\lfloor \sqrt{n} \rfloor$}
    \If{$n \bmod i = 0$}
      \State \Return \textsc{False}
    \EndIf
  \EndFor
  \State \Return \textsc{True}
\EndFunction
\end{algorithmic}
\end{algorithm}

\section{Trích dẫn tài liệu tham khảo}

Tài liệu tham khảo khai báo trong file \code{refs.bib}, mỗi tài liệu một
khối. Trong thân bài em gọi bằng lệnh \code{cite} kèm khoá của tài liệu, ví
dụ mô hình AlexNet~\cite{Krizhevsky2012} và mạng ResNet~\cite{He2016}. LaTeX
tự đánh số và tự dựng danh mục ở cuối báo cáo, nên em không phải đánh số tay.

\begin{ghinho}
Chỉ liệt kê trong \code{refs.bib} những tài liệu em thực sự trích dẫn. Tài
liệu khai báo mà không gọi \code{cite} sẽ không xuất hiện trong danh mục.
\end{ghinho}

\section{Chú thích cuối trang}

Khi cần giải thích thêm mà không muốn làm gãy mạch câu, dùng lệnh
\code{footnote}\footnote{Đây là một dòng chú thích cuối trang.}. Đừng lạm
dụng, mỗi trang nhiều lắm một đến hai chú thích.

\section{Các hộp nhấn mạnh}

Mẫu này có sẵn bốn loại hộp để làm nổi những ý quan trọng.

\begin{dinhnghia}[Học chuyển giao]
Học chuyển giao \en{(transfer learning)} là kỹ thuật tận dụng tri thức mà mô
hình đã học được trên một bài toán nguồn có nhiều dữ liệu, để giải một bài
toán đích có ít dữ liệu hơn.
\end{dinhnghia}

\begin{ghinho}
Hộp ghi nhớ dùng cho ý chốt mà em muốn người đọc mang theo sau khi đọc xong
mục đó.
\end{ghinho}

\begin{luuy}
Hộp lưu ý dùng cho cảnh báo, điều kiện áp dụng, hoặc lỗi mà người đọc dễ mắc
nếu làm theo.
\end{luuy}

\begin{vidu}[minh hoạ một tình huống cụ thể]
Hộp ví dụ tự đánh số theo chương, tiện khi em cần dẫn lại một ví dụ ở chỗ
khác trong báo cáo.
\end{vidu}


% ---------------------------------------------------------------------
%  TÀI LIỆU THAM KHẢO VÀ PHỤ LỤC
% ---------------------------------------------------------------------
\backmatter

% \phantomsection tạo mốc neo riêng cho mục lục, nếu thiếu thì hyperref
% báo trùng neo giữa dấu trang và dấu trang cha.
\cleardoublepage
\phantomsection
\addcontentsline{toc}{chapter}{Tài liệu tham khảo}
\bibliographystyle{plain}
\bibliography{refs}

% ---------------------------------------------------------------------
%  PHỤ LỤC
% ---------------------------------------------------------------------
% Trong \backmatter thì \chapter đã tự thêm vào mục lục, không cần
% \addcontentsline nữa, thêm vào là trùng neo dấu trang.
\chapter{Phụ lục}
\markboth{Phụ lục}{}

\begin{chuy}
\goiy{phụ lục là nơi đặt những nội dung dài không tiện chèn vào thân báo
cáo: mã nguồn đầy đủ, bảng dữ liệu mở rộng, kết quả thử nghiệm bổ sung, biên
bản hướng dẫn. Riêng mục A và mục B nên giữ lại.}
\end{chuy}

\section*{A. Hướng dẫn cài đặt và chạy chương trình}
\phantomsection
\addcontentsline{toc}{section}{A. Hướng dẫn cài đặt và chạy chương trình}

Viết các bước để một người chưa từng thấy đồ án vẫn chạy được sản phẩm. Hội
đồng sẽ chạy thử theo đúng các bước này.

\begin{lstlisting}[language=bash, caption={Các bước cài đặt và khởi chạy ứng dụng.}]
# 1. Clone the repository
git clone https://github.com/tencuaem/ten-do-an.git
cd ten-do-an

# 2. Create a virtual environment and install dependencies
python -m venv .venv
.venv\Scripts\activate
pip install -r requirements.txt

# 3. Train the model (skip if a trained checkpoint is provided)
python train.py --epochs 30 --data data/

# 4. Launch the web application, then open http://localhost:5000
python app.py
\end{lstlisting}

\section*{B. Bảng tự kiểm trước khi nộp}
\phantomsection
\addcontentsline{toc}{section}{B. Bảng tự kiểm trước khi nộp}

Trước khi in quyển, em tự đánh dấu vào từng ô dưới đây.

\begin{tukiem}
  \item Sản phẩm chạy được trên máy khác, không chỉ trên máy của em.
  \item Đã chạy thử demo trên đúng môi trường sẽ dùng khi bảo vệ.
  \item Mã nguồn nằm trên kho \code{git}, lịch sử commit cho thấy tiến độ
        trải đều chứ không dồn vào tuần cuối.
  \item Có tệp \code{README} và \code{requirements.txt} để người khác chạy lại.
  \item Mọi hình, bảng và giải thuật đều có chú thích, có nhãn, và được nhắc
        tới trong thân bài.
  \item Mọi mục tiêu ở Chương~\ref{ch:modau} đều được đối chiếu ở
        Chương~\ref{ch:ketluan}.
  \item Mọi con số trong báo cáo trùng khớp với kết quả chạy thật, không có
        số nào là ước lượng hay bịa.
  \item Mọi tài liệu trong \code{refs.bib} đều được trích dẫn ít nhất một
        lần, và mọi nguồn dữ liệu, mã tham khảo đều đã ghi rõ.
  \item Đã khai báo phần nào có dùng công cụ trí tuệ nhân tạo hỗ trợ, và em
        giải thích được toàn bộ hệ thống khi bị hỏi.
  \item Đã xoá hết hộp gợi ý, mọi khối nội dung mẫu, và cả
        Chương~\ref{ch:soanthao} vốn chỉ là sổ tra cứu.
\end{tukiem}

\begin{luuy}
Hai mục cuối là lỗi hay gặp nhất khi dùng mẫu có sẵn. Trước khi nộp hãy tìm
trong bản PDF xem còn dòng chữ xám \emph{NỘI DUNG MẪU} nào sót lại không.
\end{luuy}

\section*{C. Nhận xét của giảng viên hướng dẫn}
\phantomsection
\addcontentsline{toc}{section}{C. Nhận xét của giảng viên hướng dẫn}

\khungnhanxet[7.5cm]

\vspace{0.8em}
\noindent\begin{minipage}[t]{0.5\textwidth}
\color{daunavy}
{\color{daumaroon}\bfseries Điểm số}

\vspace{0.7em}
\dongchamnhan{Bằng số:}

\vspace{0.9em}
\dongchamnhan{Bằng chữ:}
\end{minipage}\hfill
\okyten{\dauGVHD}{Giảng viên hướng dẫn}


\end{document}
